% Options for packages loaded elsewhere
\PassOptionsToPackage{unicode}{hyperref}
\PassOptionsToPackage{hyphens}{url}
%
\documentclass[
  10pt,
  ignorenonframetext,
]{beamer}
\usepackage{pgfpages}
\setbeamertemplate{caption}[numbered]
\setbeamertemplate{caption label separator}{: }
\setbeamercolor{caption name}{fg=normal text.fg}
\beamertemplatenavigationsymbolsempty
% Prevent slide breaks in the middle of a paragraph
\widowpenalties 1 10000
\raggedbottom
\setbeamertemplate{part page}{
  \centering
  \begin{beamercolorbox}[sep=16pt,center]{part title}
    \usebeamerfont{part title}\insertpart\par
  \end{beamercolorbox}
}
\setbeamertemplate{section page}{
  \centering
  \begin{beamercolorbox}[sep=12pt,center]{part title}
    \usebeamerfont{section title}\insertsection\par
  \end{beamercolorbox}
}
\setbeamertemplate{subsection page}{
  \centering
  \begin{beamercolorbox}[sep=8pt,center]{part title}
    \usebeamerfont{subsection title}\insertsubsection\par
  \end{beamercolorbox}
}
\AtBeginPart{
  \frame{\partpage}
}
\AtBeginSection{
  \ifbibliography
  \else
    \frame{\sectionpage}
  \fi
}
\AtBeginSubsection{
  \frame{\subsectionpage}
}
\usepackage{amsmath,amssymb}
\usepackage{iftex}
\ifPDFTeX
  \usepackage[T1]{fontenc}
  \usepackage[utf8]{inputenc}
  \usepackage{textcomp} % provide euro and other symbols
\else % if luatex or xetex
  \usepackage{unicode-math} % this also loads fontspec
  \defaultfontfeatures{Scale=MatchLowercase}
  \defaultfontfeatures[\rmfamily]{Ligatures=TeX,Scale=1}
\fi
\usepackage{lmodern}
\usetheme[]{Berlin}
\usecolortheme{dolphin}
\usefonttheme{serif} % use mainfont rather than sansfont for slide text
\ifPDFTeX\else
  % xetex/luatex font selection
  \setmainfont[]{Calibri}
  \setsansfont[]{Calibri}
  \setmonofont[]{Consolas}
\fi
% Use upquote if available, for straight quotes in verbatim environments
\IfFileExists{upquote.sty}{\usepackage{upquote}}{}
\IfFileExists{microtype.sty}{% use microtype if available
  \usepackage[]{microtype}
  \UseMicrotypeSet[protrusion]{basicmath} % disable protrusion for tt fonts
}{}
\makeatletter
\@ifundefined{KOMAClassName}{% if non-KOMA class
  \IfFileExists{parskip.sty}{%
    \usepackage{parskip}
  }{% else
    \setlength{\parindent}{0pt}
    \setlength{\parskip}{6pt plus 2pt minus 1pt}}
}{% if KOMA class
  \KOMAoptions{parskip=half}}
\makeatother
\usepackage{xcolor}
\newif\ifbibliography
\usepackage{color}
\usepackage{fancyvrb}
\newcommand{\VerbBar}{|}
\newcommand{\VERB}{\Verb[commandchars=\\\{\}]}
\DefineVerbatimEnvironment{Highlighting}{Verbatim}{commandchars=\\\{\}}
% Add ',fontsize=\small' for more characters per line
\usepackage{framed}
\definecolor{shadecolor}{RGB}{248,248,248}
\newenvironment{Shaded}{\begin{snugshade}}{\end{snugshade}}
\newcommand{\AlertTok}[1]{\textcolor[rgb]{0.94,0.16,0.16}{#1}}
\newcommand{\AnnotationTok}[1]{\textcolor[rgb]{0.56,0.35,0.01}{\textbf{\textit{#1}}}}
\newcommand{\AttributeTok}[1]{\textcolor[rgb]{0.13,0.29,0.53}{#1}}
\newcommand{\BaseNTok}[1]{\textcolor[rgb]{0.00,0.00,0.81}{#1}}
\newcommand{\BuiltInTok}[1]{#1}
\newcommand{\CharTok}[1]{\textcolor[rgb]{0.31,0.60,0.02}{#1}}
\newcommand{\CommentTok}[1]{\textcolor[rgb]{0.56,0.35,0.01}{\textit{#1}}}
\newcommand{\CommentVarTok}[1]{\textcolor[rgb]{0.56,0.35,0.01}{\textbf{\textit{#1}}}}
\newcommand{\ConstantTok}[1]{\textcolor[rgb]{0.56,0.35,0.01}{#1}}
\newcommand{\ControlFlowTok}[1]{\textcolor[rgb]{0.13,0.29,0.53}{\textbf{#1}}}
\newcommand{\DataTypeTok}[1]{\textcolor[rgb]{0.13,0.29,0.53}{#1}}
\newcommand{\DecValTok}[1]{\textcolor[rgb]{0.00,0.00,0.81}{#1}}
\newcommand{\DocumentationTok}[1]{\textcolor[rgb]{0.56,0.35,0.01}{\textbf{\textit{#1}}}}
\newcommand{\ErrorTok}[1]{\textcolor[rgb]{0.64,0.00,0.00}{\textbf{#1}}}
\newcommand{\ExtensionTok}[1]{#1}
\newcommand{\FloatTok}[1]{\textcolor[rgb]{0.00,0.00,0.81}{#1}}
\newcommand{\FunctionTok}[1]{\textcolor[rgb]{0.13,0.29,0.53}{\textbf{#1}}}
\newcommand{\ImportTok}[1]{#1}
\newcommand{\InformationTok}[1]{\textcolor[rgb]{0.56,0.35,0.01}{\textbf{\textit{#1}}}}
\newcommand{\KeywordTok}[1]{\textcolor[rgb]{0.13,0.29,0.53}{\textbf{#1}}}
\newcommand{\NormalTok}[1]{#1}
\newcommand{\OperatorTok}[1]{\textcolor[rgb]{0.81,0.36,0.00}{\textbf{#1}}}
\newcommand{\OtherTok}[1]{\textcolor[rgb]{0.56,0.35,0.01}{#1}}
\newcommand{\PreprocessorTok}[1]{\textcolor[rgb]{0.56,0.35,0.01}{\textit{#1}}}
\newcommand{\RegionMarkerTok}[1]{#1}
\newcommand{\SpecialCharTok}[1]{\textcolor[rgb]{0.81,0.36,0.00}{\textbf{#1}}}
\newcommand{\SpecialStringTok}[1]{\textcolor[rgb]{0.31,0.60,0.02}{#1}}
\newcommand{\StringTok}[1]{\textcolor[rgb]{0.31,0.60,0.02}{#1}}
\newcommand{\VariableTok}[1]{\textcolor[rgb]{0.00,0.00,0.00}{#1}}
\newcommand{\VerbatimStringTok}[1]{\textcolor[rgb]{0.31,0.60,0.02}{#1}}
\newcommand{\WarningTok}[1]{\textcolor[rgb]{0.56,0.35,0.01}{\textbf{\textit{#1}}}}
\setlength{\emergencystretch}{3em} % prevent overfull lines
\providecommand{\tightlist}{%
  \setlength{\itemsep}{0pt}\setlength{\parskip}{0pt}}
\setcounter{secnumdepth}{-\maxdimen} % remove section numbering
% ----------------------------
% Fonts
% ----------------------------
\usepackage{fontspec}
\setmainfont{Calibri}
\setsansfont{Calibri}
\setmonofont{Consolas}
\renewcommand{\familydefault}{\sfdefault}

% ----------------------------
% Packages usuels
% ----------------------------
\usepackage{amssymb,amsmath,amsthm,enumerate}
\usepackage{array}
\usepackage{parskip}
\usepackage{graphicx}
\usepackage{caption}
\captionsetup[figure]{labelformat=empty}
\usepackage{subcaption}
\usepackage{bm}
\usepackage{multicol}
\usepackage{tcolorbox}
\usepackage{siunitx}
\usepackage{booktabs}

% ----------------------------
% Beamer navigation
% ----------------------------
% Supprimer les icônes en bas
\setbeamertemplate{navigation symbols}{}

% Activer les mini-frames (petits ronds)
\useoutertheme{miniframes}

% Section en gras dans la barre du haut
\setbeamerfont{section in head/foot}{series=\bfseries}

% ----------------------------
% Headline sur 2 lignes :
% 1) Sections (en haut)
% 2) Mini-frames seuls (en dessous)
% ----------------------------
\makeatletter
\setbeamertemplate{headline}{%
  % --- Ligne 1 : sections
  \begin{beamercolorbox}[wd=\paperwidth,ht=2.8ex,dp=1.2ex]{section in head/foot}%
    \insertsectionnavigationhorizontal{\paperwidth}{}{}%
  \end{beamercolorbox}%
}
\makeatother

% ----------------------------
% (Optionnel) : pas de barre sur la slide de titre
% ----------------------------
\addtobeamertemplate{title page}{\setbeamertemplate{headline}{}}{}
\ifLuaTeX
  \usepackage{selnolig}  % disable illegal ligatures
\fi
\usepackage{bookmark}
\IfFileExists{xurl.sty}{\usepackage{xurl}}{} % add URL line breaks if available
\urlstyle{same}
\hypersetup{
  pdftitle={R pour jeunes économistes},
  pdfauthor={J.-B. Guiffard},
  hidelinks,
  pdfcreator={LaTeX via pandoc}}

\title{R pour jeunes économistes}
\subtitle{RMarkdown et production de documents}
\author{J.-B. Guiffard}
\date{5 juin 2026}

\begin{document}
\frame{\titlepage}

\section{RMarkdown et production de
documents}\label{rmarkdown-et-production-de-documents}

\begin{frame}{Objet de la séance}
\phantomsection\label{objet-de-la-suxe9ance}
Cette dernière séance est consacrée à la production de documents avec
\textbf{RMarkdown}.

L'objectif est de comprendre comment passer :

\[
\text{script R} \rightarrow \text{document reproductible} \rightarrow \text{rapport / slides / html / pdf}
\]

Nous verrons comment utiliser RMarkdown pour :

\begin{itemize}
\tightlist
\item
  rédiger du texte ;
\item
  intégrer du code R ;
\item
  produire automatiquement des tableaux ;
\item
  intégrer des graphiques ;
\item
  générer plusieurs formats de sortie ;
\item
  automatiser la production de rapports.
\end{itemize}

\(\Rightarrow\) RMarkdown permet de réunir dans un même fichier le
raisonnement, le code et les résultats.
\end{frame}

\begin{frame}{Pourquoi utiliser RMarkdown ?}
\phantomsection\label{pourquoi-utiliser-rmarkdown}
Dans un travail de recherche, on produit souvent :

\begin{itemize}
\tightlist
\item
  des statistiques descriptives ;
\item
  des graphiques ;
\item
  des tableaux de régression ;
\item
  des cartes ;
\item
  des rapports intermédiaires ;
\item
  des supports de présentation.
\end{itemize}

Avec un workflow classique, ces éléments sont souvent produits
séparément puis copiés-collés dans Word, PowerPoint ou LaTeX.

Avec RMarkdown, on peut produire ces éléments directement à partir du
code.

\(\Rightarrow\) Cela réduit les erreurs manuelles et facilite la mise à
jour des résultats.
\end{frame}

\begin{frame}{RMarkdown comme outil de recherche reproductible}
\phantomsection\label{rmarkdown-comme-outil-de-recherche-reproductible}
RMarkdown permet de produire un document à partir d'un fichier source
unique.

Ce fichier contient :

\begin{itemize}
\tightlist
\item
  le texte ;
\item
  le code ;
\item
  les résultats ;
\item
  les graphiques ;
\item
  les tableaux ;
\item
  les options du document.
\end{itemize}

\[
\text{texte + code + données}
\rightarrow
\text{document final}
\]
\end{frame}

\begin{frame}[fragile]{Plan de la séance}
\phantomsection\label{plan-de-la-suxe9ance}
La séance est organisée en six étapes.

\begin{enumerate}
\tightlist
\item
  Comprendre le principe de RMarkdown
\item
  Créer un premier document \texttt{.Rmd}
\item
  Écrire du texte avec la syntaxe Markdown
\item
  Insérer du code R et contrôler les chunks
\item
  Intégrer tableaux, graphiques et résultats
\item
  Automatiser et paramétrer des rapports
\end{enumerate}
\end{frame}

\section{Comprendre RMarkdown}\label{comprendre-rmarkdown}

\begin{frame}[fragile]{Qu'est-ce qu'un fichier RMarkdown ?}
\phantomsection\label{quest-ce-quun-fichier-rmarkdown}
Un fichier RMarkdown est un fichier texte avec l'extension
\texttt{.Rmd}.

Il permet de mélanger :

\begin{itemize}
\tightlist
\item
  du texte ;
\item
  des titres ;
\item
  des équations ;
\item
  du code R ;
\item
  des sorties produites par R.
\end{itemize}

Exemple :

mean(c(1, 2, 3, 4, 5))

Après compilation, le résultat est intégré dans le document final.
\end{frame}

\begin{frame}[fragile]{Le principe général}
\phantomsection\label{le-principe-guxe9nuxe9ral}
RMarkdown fonctionne en deux étapes.

\[
\text{.Rmd}
\rightarrow
\text{.md}
\rightarrow
\text{document final}
\]

Première étape :

\begin{itemize}
\tightlist
\item
  le package \texttt{knitr} exécute le code R ;
\item
  les résultats sont insérés dans un document Markdown.
\end{itemize}

Deuxième étape :

\begin{itemize}
\tightlist
\item
  \texttt{pandoc} transforme le document en HTML, PDF, Word ou slides.
\end{itemize}
\end{frame}

\begin{frame}[fragile]{Quels formats peut-on produire ?}
\phantomsection\label{quels-formats-peut-on-produire}
À partir d'un même fichier \texttt{.Rmd}, on peut produire plusieurs
formats.

\begin{itemize}
  \item \texttt{html\_document} : rapport HTML ;
  \item \texttt{pdf\_document} : rapport PDF ;
  \item \texttt{word\_document} : document Word ;
  \item \texttt{beamer\_presentation} : slides PDF ;
  \item \texttt{ioslides\_presentation} : slides HTML.
\end{itemize}

Le contenu reste le même, mais le format de sortie peut changer.
\end{frame}

\begin{frame}{Pourquoi c'est utile pour un doctorant ?}
\phantomsection\label{pourquoi-cest-utile-pour-un-doctorant}
RMarkdown est particulièrement utile pour :

\begin{itemize}
\tightlist
\item
  produire des rapports intermédiaires ;
\item
  documenter une analyse empirique ;
\item
  préparer un support de séminaire ;
\item
  partager des résultats avec des co-auteurs ;
\item
  produire des annexes reproductibles ;
\item
  automatiser des tableaux ou graphiques.
\end{itemize}

Exemple : Lorsque les données changent, il suffit de recompiler le
document pour mettre à jour les résultats.
\end{frame}

\section{Créer un premier document
RMarkdown}\label{cruxe9er-un-premier-document-rmarkdown}

\begin{frame}{Structure générale d'un document}
\phantomsection\label{structure-guxe9nuxe9rale-dun-document}
Un document RMarkdown contient deux grandes parties :

\begin{enumerate}
\tightlist
\item
  un en-tête YAML ;
\item
  un corps de document.
\end{enumerate}

L'en-tête définit les métadonnées et le format de sortie.

Le corps contient le texte, le code et les résultats.

title: ``Analyse des émissions de CO2'' author: ``Prénom Nom'' date:
``2026-06-05'' output: html\_document
\end{frame}

\begin{frame}{Exemple d'en-tête simple}
\phantomsection\label{exemple-den-tuxeate-simple}
title: ``Analyse des émissions de CO2'' subtitle: ``Exemple de rapport
reproductible'' author: ``Jean-Baptiste Guiffard'' date: ``5 juin 2026''
output: html\_document

Cet en-tête indique :

\begin{itemize}
\tightlist
\item
  le titre du document ;
\item
  le sous-titre ;
\item
  l'auteur ;
\item
  la date ;
\item
  le format de sortie.
\end{itemize}
\end{frame}

\begin{frame}{Plusieurs formats de sortie}
\phantomsection\label{plusieurs-formats-de-sortie}
On peut demander plusieurs formats dans le même fichier.

title: ``Analyse des émissions de CO2'' author: ``Prénom Nom'' output:
html\_document: toc: true pdf\_document: toc: true word\_document:
default

Cela permet de produire :

\begin{itemize}
\tightlist
\item
  un rapport HTML ;
\item
  un rapport PDF ;
\item
  un document Word.
\end{itemize}
\end{frame}

\begin{frame}[fragile]{Le chunk setup}
\phantomsection\label{le-chunk-setup}
Au début du document, on ajoute souvent un chunk appelé \texttt{setup}.

knitr::opts\_chunk\$set( echo = TRUE, warning = FALSE, message = FALSE )

Ce chunk permet de définir les options par défaut.

Ici :

\begin{itemize}
\tightlist
\item
  le code est affiché ;
\item
  les warnings ne sont pas affichés ;
\item
  les messages ne sont pas affichés.
\end{itemize}
\end{frame}

\section{Écrire du texte avec
Markdown}\label{uxe9crire-du-texte-avec-markdown}

\begin{frame}[fragile]{Titres, texte et listes}
\phantomsection\label{titres-texte-et-listes}
La syntaxe Markdown est légère.

\begin{Shaded}
\begin{Highlighting}[]
\FunctionTok{\# Titre de niveau 1}

\FunctionTok{\#\# Titre de niveau 2}

\FunctionTok{\#\#\# Titre de niveau 3}

\NormalTok{Ceci est du texte avec *de l’italique* et **du gras**.}

\SpecialStringTok{{-} }\NormalTok{premier élément}
\SpecialStringTok{{-} }\NormalTok{deuxième élément}
\SpecialStringTok{{-} }\NormalTok{troisième élément}
\end{Highlighting}
\end{Shaded}

Markdown permet d'écrire un document lisible même avant compilation.
\end{frame}

\begin{frame}[fragile]{Liens, images et équations}
\phantomsection\label{liens-images-et-uxe9quations}
On peut insérer un lien :

\begin{Shaded}
\begin{Highlighting}[]
\CommentTok{[}\OtherTok{Our World in Data}\CommentTok{](https://ourworldindata.org/)}
\end{Highlighting}
\end{Shaded}

On peut insérer une image :

\begin{Shaded}
\begin{Highlighting}[]
\AlertTok{![Titre de la figure](figures/ma\_figure.png)}\NormalTok{\{width=70\%\}}
\end{Highlighting}
\end{Shaded}

On peut écrire une équation :

\begin{Shaded}
\begin{Highlighting}[]
\NormalTok{$$}
\NormalTok{Y\_i = \textbackslash{}alpha + \textbackslash{}beta X\_i + \textbackslash{}varepsilon\_i}
\NormalTok{$$}
\end{Highlighting}
\end{Shaded}
\end{frame}

\begin{frame}{Insérer un résultat dans le texte}
\phantomsection\label{insuxe9rer-un-ruxe9sultat-dans-le-texte}
On peut insérer un résultat R directement dans une phrase.

La base contient r nrow(co2\_analysis) observations.

Après compilation, R remplace le code par le résultat.

Exemple :

\begin{quote}
La base contient 4231 observations.
\end{quote}

Le code inline évite de recopier manuellement des chiffres dans le
texte.
\end{frame}

\section{Insérer du code R}\label{insuxe9rer-du-code-r}

\begin{frame}{Qu'est-ce qu'un chunk ?}
\phantomsection\label{quest-ce-quun-chunk}
Un chunk est un bloc de code R inséré dans le document.

\{r\}

Lors de la compilation :

\begin{itemize}
\tightlist
\item
  le code peut être exécuté ;
\item
  le résultat peut être affiché ;
\item
  les graphiques peuvent être intégrés.
\end{itemize}
\end{frame}

\begin{frame}{Charger les packages}
\phantomsection\label{charger-les-packages}
Dans un rapport reproductible, les packages doivent être chargés dans le
document.

\{r packages, message=FALSE, warning=FALSE\}

Cela permet au document de fonctionner dans une nouvelle session R.
\end{frame}

\begin{frame}{Importer les données}
\phantomsection\label{importer-les-donnuxe9es}
Les données doivent aussi être importées dans le document.

\{r import-data, eval=FALSE\} co2\_raw \textless-
read\_csv(``data/raw/owid-co2-data.csv'')

On utilise de préférence des chemins relatifs.

À éviter :

\{r, eval=FALSE\}
read\_csv(``C:/Users/Jean-Baptiste/Desktop/owid-co2-data.csv'')
\end{frame}

\begin{frame}{Construire une base d'analyse}
\phantomsection\label{construire-une-base-danalyse}
On peut construire la base directement dans le document.

\{r build-data, eval=FALSE\} co2\_analysis \textless- co2\_raw
\%\textgreater\% filter( year \textgreater= 1990, !is.na(iso\_code),
!is.na(co2\_per\_capita), !is.na(population), !is.na(gdp), population
\textgreater{} 0, gdp \textgreater{} 0, co2\_per\_capita \textgreater{}
0 ) \%\textgreater\% mutate( gdp\_per\_capita = gdp / population,
log\_gdp\_pc = log(gdp\_per\_capita), log\_co2\_pc =
log(co2\_per\_capita) )
\end{frame}

\section{Contrôler les chunks}\label{contruxf4ler-les-chunks}

\begin{frame}{Les principales options}
\phantomsection\label{les-principales-options}
Les chunks peuvent être contrôlés avec des options.

\begin{itemize}
  \item \texttt{echo} : afficher ou non le code ;
  \item \texttt{eval} : exécuter ou non le code ;
  \item \texttt{include} : inclure ou non le code et les résultats ;
  \item \texttt{warning} : afficher ou non les warnings ;
  \item \texttt{message} : afficher ou non les messages ;
  \item \texttt{fig.width} et \texttt{fig.height} : taille des figures.
\end{itemize}
\end{frame}

\begin{frame}{Exemple : afficher le résultat mais pas le code}
\phantomsection\label{exemple-afficher-le-ruxe9sultat-mais-pas-le-code}
echo=FALSE

\{r desc-table, eval=FALSE\} co2\_analysis \%\textgreater\% summarise(
observations = n(), countries = n\_distinct(country), mean\_co2\_pc =
mean(co2\_per\_capita, na.rm = TRUE) )

\begin{itemize}
\tightlist
\item
  le code est exécuté ;
\item
  le résultat est affiché ;
\item
  le code n'apparaît pas dans le document.
\end{itemize}
\end{frame}

\begin{frame}{Exemple : montrer le code sans l'exécuter}
\phantomsection\label{exemple-montrer-le-code-sans-lexuxe9cuter}
\{r exemple-code, eval=FALSE\} co2\_analysis \%\textgreater\%
group\_by(year) \%\textgreater\% summarise(mean\_co2 =
mean(co2\_per\_capita, na.rm = TRUE))

\begin{itemize}
\tightlist
\item
  le code est affiché ;
\item
  le code n'est pas exécuté ;
\item
  aucun résultat n'est produit.
\end{itemize}
\end{frame}

\section{Tableaux dans RMarkdown}\label{tableaux-dans-rmarkdown}

\begin{frame}[fragile]{Produire une table descriptive}
\phantomsection\label{produire-une-table-descriptive}
On peut produire une table à partir du code R.

\{r, eval=FALSE\}

desc\_table \textless- co2\_analysis \%\textgreater\% summarise(
Observations = n(), Pays = n\_distinct(country),
\texttt{CO2\ moyen\ par\ habitant} = mean(co2\_per\_capita, na.rm =
TRUE), \texttt{PIB\ moyen\ par\ habitant} = mean(gdp\_per\_capita, na.rm
= TRUE) )

knitr::kable(desc\_table, digits = 2)

Les chiffres du tableau sont recalculés automatiquement à chaque
compilation.
\end{frame}

\begin{frame}{Table par groupe}
\phantomsection\label{table-par-groupe}
On peut produire des tableaux par année ou par pays.

\{r, eval=FALSE\}

year\_stats \textless- co2\_analysis \%\textgreater\% group\_by(year)
\%\textgreater\% summarise( mean\_co2\_pc = mean(co2\_per\_capita, na.rm
= TRUE), median\_co2\_pc = median(co2\_per\_capita, na.rm = TRUE),
n\_countries = n\_distinct(country), .groups = ``drop'' )

knitr::kable(head(year\_stats, 10), digits = 2)
\end{frame}

\begin{frame}{Tableaux plus avancés}
\phantomsection\label{tableaux-plus-avancuxe9s}
Pour des besoins plus avancés, on peut utiliser plusieurs packages.

\begin{itemize}
  \item \texttt{knitr::kable()} : tableaux simples ;
  \item \texttt{kableExtra} : mise en forme avancée ;
  \item \texttt{gt} : tableaux modernes ;
  \item \texttt{flextable} : tableaux pour Word et PowerPoint ;
  \item \texttt{modelsummary} : tableaux de régressions et statistiques descriptives.
\end{itemize}

\begin{block}{Conseil}
Commencer simple avec \texttt{kable()}, puis utiliser des packages plus avancés selon les besoins.
\end{block}
\end{frame}

\section{Graphiques dans RMarkdown}\label{graphiques-dans-rmarkdown}

\begin{frame}{Intégrer un graphique ggplot}
\phantomsection\label{intuxe9grer-un-graphique-ggplot}
Les graphiques produits par R sont intégrés automatiquement.

\{r graph-distribution, fig.width=7, fig.height=4\}

\{r, eval=FALSE\} ggplot(co2\_analysis, aes(x = log\_co2\_pc)) +
geom\_histogram(bins = 40) + theme\_minimal() + labs( title =
``Distribution du log des émissions de CO2 par habitant'', x = ``log(CO2
par habitant)'', y = ``Nombre d'observations'' )
\end{frame}

\begin{frame}{Graphique temporel}
\phantomsection\label{graphique-temporel}
\{r graph-time, fig.width=7, fig.height=4\} \{r, eval=FALSE\}

co2\_year \textless- co2\_analysis \%\textgreater\% group\_by(year)
\%\textgreater\% summarise( mean\_co2\_pc = mean(co2\_per\_capita, na.rm
= TRUE), .groups = ``drop'' )

ggplot(co2\_year, aes(x = year, y = mean\_co2\_pc)) + geom\_line() +
theme\_minimal() + labs( title = ``Évolution moyenne des émissions de
CO2 par habitant'', x = ``Année'', y = ``CO2 par habitant moyen'' )

Nommer les chunks facilite la résolution des erreurs.
\end{frame}

\section{Régressions et résultats}\label{ruxe9gressions-et-ruxe9sultats}

\begin{frame}{Une première régression}
\phantomsection\label{une-premiuxe8re-ruxe9gression}
On peut estimer un modèle directement dans le document.

\{r, eval=FALSE\} model\_1 \textless- lm(log\_co2\_pc \textasciitilde{}
log\_gdp\_pc, data = co2\_analysis)

summary(model\_1)

Cette régression décrit une relation entre PIB par habitant et émissions
de CO2 par habitant.
\end{frame}

\begin{frame}[fragile]{Présenter une table de régression}
\phantomsection\label{pruxe9senter-une-table-de-ruxe9gression}
\{r regression-table, echo=FALSE\}

\{r, eval=FALSE\}

modelsummary::modelsummary( model\_1, stars = TRUE, title = ``Relation
entre PIB par habitant et émissions de CO2'' )

Le tableau est produit automatiquement à partir de l'objet
\texttt{model\_1}.
\end{frame}

\section{Rapports automatisés}\label{rapports-automatisuxe9s}

\begin{frame}{Pourquoi automatiser un rapport ?}
\phantomsection\label{pourquoi-automatiser-un-rapport}
Un rapport automatisé est utile lorsqu'un document doit être mis à jour
régulièrement.

Exemples :

\begin{itemize}
\tightlist
\item
  rapport mensuel ;
\item
  tableau de bord ;
\item
  note de suivi ;
\item
  rapport par pays ;
\item
  rapport par région ;
\item
  rapport par année ;
\item
  annexe empirique.
\end{itemize}

Le rapport est produit à partir du code : si les données changent, le
document peut être mis à jour automatiquement.
\end{frame}

\begin{frame}{Compiler avec render()}
\phantomsection\label{compiler-avec-render}
On peut compiler un fichier RMarkdown depuis R.

\{r, eval=FALSE\} rmarkdown::render(``rapport\_co2.Rmd'')

On peut aussi choisir un environnement propre.

\{r, eval=FALSE\} rmarkdown::render( input = ``rapport\_co2.Rmd'', envir
= new.env() )
\end{frame}

\begin{frame}{Rapports paramétrés}
\phantomsection\label{rapports-paramuxe9truxe9s}
Un rapport paramétré permet de produire plusieurs documents à partir du
même modèle.

Exemple : produire un rapport pour un pays donné.

Dans le YAML :

title: ``Rapport CO2'' output: html\_document params: country\_name:
``France''

Dans le code :

\{r, eval=FALSE\} co2\_country \textless- co2\_analysis \%\textgreater\%
filter(country == params\$country\_name)
\end{frame}

\begin{frame}{Générer un rapport pour un pays}
\phantomsection\label{guxe9nuxe9rer-un-rapport-pour-un-pays}
\{r, eval=FALSE\} rmarkdown::render( input =
``rapport\_co2\_parametre.Rmd'', params = list(country\_name =
``Brazil''), output\_file = ``rapport\_co2\_Brazil.html'', envir =
new.env() )

On peut produire le même rapport pour un autre pays :

\{r, eval=FALSE\} rmarkdown::render( input =
``rapport\_co2\_parametre.Rmd'', params = list(country\_name =
``India''), output\_file = ``rapport\_co2\_India.html'', envir =
new.env() )
\end{frame}

\begin{frame}{Générer plusieurs rapports}
\phantomsection\label{guxe9nuxe9rer-plusieurs-rapports}
\{r, eval=FALSE\} countries \textless- c(``France'', ``Brazil'',
``India'')

for (pays in countries) \{ rmarkdown::render( input =
``rapport\_co2\_parametre.Rmd'', params = list(country\_name = pays),
output\_file = paste0(``rapport\_co2\_'', pays, ``.html''), envir =
new.env() ) \}

\begin{block}{Usage}
Très utile pour produire automatiquement des rapports par pays, région, secteur ou année.
\end{block}
\end{frame}

\section{Bonnes pratiques}\label{bonnes-pratiques}

\begin{frame}[fragile]{Organiser son projet}
\phantomsection\label{organiser-son-projet}
Un projet RMarkdown doit être organisé clairement.

\begin{Shaded}
\begin{Highlighting}[]
\NormalTok{projet\_rmarkdown/}
\NormalTok{├── data/}
\NormalTok{│   ├── raw/}
\NormalTok{│   └── clean/}
\NormalTok{├── scripts/}
\NormalTok{├── reports/}
\NormalTok{├── outputs/}
\NormalTok{│   ├── figures/}
\NormalTok{│   └── tables/}
\NormalTok{└── projet\_rmarkdown.Rproj}
\end{Highlighting}
\end{Shaded}
\end{frame}

\begin{frame}{Les erreurs classiques}
\phantomsection\label{les-erreurs-classiques}
Quelques erreurs reviennent souvent avec RMarkdown.

\begin{itemize}
  \item Utiliser des chemins absolus.
  \item Oublier de charger un package dans le document.
  \item Utiliser un objet présent dans l’environnement mais non créé dans le `.Rmd`.
  \item Ne pas nommer les chunks.
  \item Copier-coller manuellement des chiffres dans le texte.
  \item Masquer les messages d’erreur au lieu de les comprendre.
  \item Compiler rarement le document.
  \item Ne pas vérifier que le document compile dans une nouvelle session.
\end{itemize}
\end{frame}

\begin{frame}{Conclusion}
\phantomsection\label{conclusion}
\centering
\Large

Merci beaucoup pour votre attention !

\vspace{0.5cm}

\large

Dernière discussion :

Comment pourriez-vous utiliser RMarkdown dans votre propre thèse ?
\end{frame}

\end{document}
